\chapter{Results}
\label{chap:results}

\section{Database coverage and imbalance}
\label{sec:database-results}

\subsection{Scale, source concentration, and target availability}

Before task-specific cleaning, the assembled table contains 5,057 experimental
rows mapped to 87 publications. The raw database overview in
\cref{fig:database-overview} shows its breadth: pressure spans more than three
orders of magnitude, gravity extends from microgravity to hypergravity, and
2,590 rows contain a numeric FSR before the regression pipeline's stricter
parsing and deduplication. After task-specific preprocessing, 4,487 rows remain
for ignition classification and 2,531 for FSR regression.

\begin{figure}[htbp]
  \centering
  \includegraphics[width=0.88\textwidth]{database_overview.png}
  \caption{Raw assembled-database overview before task-specific deduplication.
  The modeling cohorts are smaller because their loaders impose distinct target,
  parsing, and duplicate rules. Source:
  \texttt{metrics/01\_dataset\_overview\_summary.png}.}
  \label{fig:database-overview}
\end{figure}

The number of rows is not the number of independent experiments. Publications
contribute very different numbers of nearby conditions. The Pareto curve in
\cref{fig:database-imbalance} shows that 36 papers contribute more than 80\% of
all rows. The ignition target is also imbalanced toward positive outcomes.
These two imbalances motivate paper weighting, prevalence-aware classification
metrics, and paper-grouped evaluation.

\begin{figure}[htbp]
  \centering
  \begin{subfigure}[t]{0.61\textwidth}
    \includegraphics[width=\linewidth]{database_pareto.png}
    \caption{Cumulative data share after sorting papers by contribution.}
  \end{subfigure}\hfill
  \begin{subfigure}[t]{0.36\textwidth}
    \includegraphics[width=\linewidth]{database_ignition_balance.png}
    \caption{Raw ignition outcome composition.}
  \end{subfigure}
  \caption{Source-size and class imbalance in the assembled database. Source:
  repository metrics figures 05 and 06.}
  \label{fig:database-imbalance}
\end{figure}

\subsection{Oxygen--pressure coverage}

\Cref{fig:o2-pressure-coverage} makes the sampling geometry visible. Most rows
cluster near terrestrial pressure and a limited range of oxygen fractions;
other pressure--oxygen combinations are represented by isolated campaigns or
not sampled at all. Ignition and no-ignition outcomes overlap substantially,
showing that neither oxygen nor pressure alone determines the label. The dense
bands and empty cells also explain why a random row split can appear much easier
than transfer to a new paper: nearby points commonly share a source.

\begin{figure}[htbp]
  \centering
  \begin{subfigure}[t]{0.49\textwidth}
    \includegraphics[width=\linewidth]{database_o2_pressure_ignition.png}
    \caption{Observed outcomes in oxygen--pressure space.}
  \end{subfigure}\hfill
  \begin{subfigure}[t]{0.49\textwidth}
    \includegraphics[width=\linewidth]{database_o2_pressure_density.png}
    \caption{Sampling density; blank cells are empirical gaps.}
  \end{subfigure}
  \caption{Database coverage over oxygen mole fraction and pressure. The
  pressure axis is logarithmic. Source: repository metrics figures 25 and 28.}
  \label{fig:o2-pressure-coverage}
\end{figure}

Gravity coverage is similarly uneven (\cref{fig:gravity-coverage}). Ground,
microgravity, lunar, Martian, and other partial-gravity regimes are present, but
their row counts and associated materials differ. A numeric gravity field does
not turn sparse partial-gravity combinations into interpolation evidence.

\begin{figure}[htbp]
  \centering
  \includegraphics[width=0.72\textwidth]{database_gravity_regimes.png}
  \caption{Gravity-regime composition in the assembled database. Source:
  \texttt{metrics/23\_gravity\_regime\_composition.png}.}
  \label{fig:gravity-coverage}
\end{figure}

\section{Ignition classification}
\label{sec:results-ignition}

\subsection{Model families and interpolation}

The selected interpolation classifier is
\texttt{xgb\_paper\_sqrt}, an XGBoost model with paper weight
\(1/\sqrt{n_p}\). Across 15 repeated stratified outer folds it reaches
ROC--AUC \(0.9076\pm0.0102\), PR--AUC \(0.9687\pm0.0039\), and mean Brier
score 0.1063. The unweighted XGBoost baseline is close
(ROC--AUC \(0.9058\pm0.0091\)); SVM, decision tree, and KNN are lower.

\Cref{fig:ignition-family-comparison} summarizes the best persisted result in
each family under both protocols. The blank MLP panel reflects the discrepancy
between available MLP shards and the formal merged selection pool; it is not a
zero score. The plot makes the principal result clearer than a long leaderboard:
every represented family performs worse when complete papers are withheld.

\begin{figure}[htbp]
  \centering
  \includegraphics[width=\textwidth]{ignition_family_comparison.png}
  \caption{Best persisted ignition candidate by family for row interpolation
  and paper-grouped extrapolation (mean ROC--AUC with fold SD). Source:
  \texttt{Ignition Classifiers/results/report/five\_model\_family\_comparison.png}.}
  \label{fig:ignition-family-comparison}
\end{figure}

\subsection{Extrapolation and generalization loss}

The extrapolation champion changes to the unweighted
\texttt{xgb\_baseline}. Its grouped ROC--AUC is
\(0.7183\pm0.0444\), PR--AUC \(0.8767\pm0.0298\), and Brier score 0.1791.
The paper-cluster bootstrap estimate is 0.7101 ROC--AUC with 95\% interval
0.6408--0.7848. LOPO robustness is similar (0.7020,
0.6376--0.7751).

For the same unweighted XGBoost recipe, ROC--AUC falls by 0.1875 between row
and grouped evaluation. \Cref{fig:ignition-generalization} shows that this
loss is broad across candidate families. Small bars for focal variants do not
indicate successful transfer; they result from near-chance performance in both
protocols.

\begin{figure}[htbp]
  \centering
  \includegraphics[width=0.80\textwidth]{ignition_generalization_drop.png}
  \caption{Ignition ROC--AUC loss when evaluation changes from random rows to
  held-out papers. Source:
  \texttt{Ignition Classifiers/results/report/interpolation\_extrapolation\_drop.png}.}
  \label{fig:ignition-generalization}
\end{figure}

\subsection{Thresholds and calibration}

Probability ranking and operational decisions are different. In grouped
evaluation, the F1-oriented threshold (approximately 0.285) gives sensitivity
0.984 but specificity only 0.059, effectively labeling almost every case as
ignition. The balanced-accuracy/Youden threshold (approximately 0.858) gives
sensitivity 0.664 and specificity 0.662. \Cref{fig:ignition-operating} shows
these trade-offs together with calibration. No threshold is universally
correct: a screening policy must state the relative cost of missed ignitions and
unnecessary tests.

\begin{figure}[htbp]
  \centering
  \begin{subfigure}[t]{0.49\textwidth}
    \includegraphics[width=\linewidth]{ignition_thresholds.png}
    \caption{Threshold-dependent champion performance.}
  \end{subfigure}\hfill
  \begin{subfigure}[t]{0.49\textwidth}
    \includegraphics[width=\linewidth]{ignition_calibration.png}
    \caption{Observed frequency versus predicted probability.}
  \end{subfigure}
  \caption{Ignition operating-point and calibration diagnostics from persisted
  outer predictions. Source: ignition report figures.}
  \label{fig:ignition-operating}
\end{figure}

\subsection{Ignition explainability}

The held-paper XGBoost SHAP plot in \cref{fig:ignition-explainability} shows
both contribution direction and feature value. Oxygen has the largest global
SHAP magnitude, followed by absolute flow speed, Prandtl number, gas molar
mass, fuel pyrolysis temperature, and ignition time. Higher oxygen values
generally move predictions toward ignition in this fitted model, while several
transport quantities show nonlinear or mixed patterns.

Permutation importance gives a different ranking: fuel pyrolysis temperature,
ignition time, fuel heat capacity, and thermal diffusion time produce the
largest positive held-out ROC--AUC losses when shuffled. Oxygen has high SHAP
magnitude but slightly negative permutation importance in this held-out sample.
The disagreement is evidence of correlated and substitutable predictors, not a
contradiction and not a causal ranking.

\begin{figure}[htbp]
  \centering
  \begin{subfigure}[t]{0.49\textwidth}
    \includegraphics[width=\linewidth]{ignition_shap_beeswarm.png}
    \caption{SHAP values for the extrapolation champion.}
  \end{subfigure}\hfill
  \begin{subfigure}[t]{0.49\textwidth}
    \includegraphics[width=\linewidth]{ignition_permutation.png}
    \caption{Held-out permutation importance.}
  \end{subfigure}
  \caption{Authoritative ignition explainability outputs produced by the
  classification pipeline. Source:
  \texttt{Ignition Classifiers/results/report/extrapolation\_champion/}.}
  \label{fig:ignition-explainability}
\end{figure}

\section{Flame-spread-rate regression}
\label{sec:results-fsr}

\subsection{Model comparison and interpolation}

The unaugmented random-forest baseline is selected for FSR interpolation. Across
ten repeated row holdouts it reaches
\(R^2=0.8790\pm0.0178\),
RMSE \(0.006263\pm0.000472~\si{\metre\per\second}\), and
MAE \(0.002103\pm0.000173~\si{\metre\per\second}\).
Bootstrap-augmented XGBoost has a numerically slightly lower mean RMSE
(\(0.006236~\si{\metre\per\second}\)) but falls within the 2\% tie band and has
greater fold variability, so the declared tie policy favors the unaugmented
random forest.

\Cref{fig:fsr-model-comparison} presents the complete persisted comparison
visually. Tree ensembles occupy the leading group. More importantly, the
orange grouped bars are consistently higher than the blue interpolation bars;
augmentation does not remove that protocol gap.

\begin{figure}[htbp]
  \centering
  \includegraphics[width=\textwidth]{fsr_model_comparison.png}
  \caption{FSR outer-fold RMSE for persisted candidates under row
  interpolation, paper extrapolation, and real/bootstrapped training conditions.
  Source:
  \texttt{FSR Regressors/results/report/01\_comparison/model\_comparison\_bars.png}.}
  \label{fig:fsr-model-comparison}
\end{figure}

\subsection{Extrapolation and paper heterogeneity}

The same unaugmented random-forest recipe is selected for paper extrapolation,
with a separately tuned configuration. Grouped performance is
\(R^2=0.4576\pm0.2215\),
RMSE \(0.012205\pm0.004115~\si{\metre\per\second}\), and
MAE \(0.004819\pm0.001144~\si{\metre\per\second}\). The RMSE is 1.95 times
the interpolation value. The paper-cluster 95\% interval for RMSE is
0.006525--\(0.017839~\si{\metre\per\second}\), demonstrating sensitivity to
which campaigns are held out.

\begin{figure}[htbp]
  \centering
  \begin{subfigure}[t]{0.49\textwidth}
    \includegraphics[width=\linewidth]{fsr_generalization_gap.png}
    \caption{Grouped minus row-level RMSE by candidate.}
  \end{subfigure}\hfill
  \begin{subfigure}[t]{0.49\textwidth}
    \includegraphics[width=\linewidth]{fsr_group_rmse_strip.png}
    \caption{Grouped RMSE across repeated paper splits.}
  \end{subfigure}
  \caption{FSR transfer loss and split-to-split heterogeneity. Error bars in
  (a) summarize repeat variability. Source: FSR report comparison figures.}
  \label{fig:fsr-generalization}
\end{figure}

The parity diagnostic in \cref{fig:fsr-parity} shows dense low-FSR observations
near the diagonal but larger errors and visible regression toward the mean at
sparse high spread rates. The hardest held-out paper reaches RMSE
\(0.1331~\si{\metre\per\second}\), whereas many papers are below
\(0.003~\si{\metre\per\second}\). A single pooled score therefore conceals
important campaign-level failure modes.

\begin{figure}[htbp]
  \centering
  \includegraphics[width=0.76\textwidth]{fsr_pred_vs_true.png}
  \caption{Observed versus predicted FSR for the random-forest family under
  paper-grouped evaluation. Source:
  \texttt{FSR Regressors/results/report/02\_diagnostics/pred\_vs\_true\_random\_forest.png}.}
  \label{fig:fsr-parity}
\end{figure}

\subsection{FSR explainability}

The FSR random forest relies most strongly on thermal diffusion time, oxygen,
gas thermal conductivity, and flow velocity in its SHAP summary
(\cref{fig:fsr-explainability}). The beeswarm adds direction: high oxygen
generally pushes FSR upward, while thermal diffusion time has a strongly
nonlinear pattern. Held-out permutation importance independently identifies
oxygen, flow velocity, gas conductivity, and thermal diffusion time as the
largest positive RMSE degradations.

Several correlated material variables have negative permutation importance in
the realized held-out sample. That means shuffling happened to reduce error; it
does not imply a physically protective effect. The repeated appearance of
oxygen, flow, and thermal transport across two explanation methods supports
their use as hypotheses and coverage priorities, not as causal coefficients.

\begin{figure}[htbp]
  \centering
  \begin{subfigure}[t]{0.49\textwidth}
    \includegraphics[width=\linewidth]{fsr_shap_beeswarm.png}
    \caption{Random-forest SHAP values on unseen papers.}
  \end{subfigure}\hfill
  \begin{subfigure}[t]{0.49\textwidth}
    \includegraphics[width=\linewidth]{fsr_permutation.png}
    \caption{Held-out permutation importance.}
  \end{subfigure}
  \caption{Authoritative FSR explainability outputs generated by the regression
  pipeline. Source:
  \texttt{FSR Regressors/results/report/07\_explainability/}.}
  \label{fig:fsr-explainability}
\end{figure}

\section{Synthesis}

The figures establish one common result across two targets and different model
families: within-paper interpolation is strong, while transfer to an unseen
paper is moderate and substantially more uncertain. Ignition ROC--AUC falls
from approximately 0.91 to 0.72; FSR RMSE approximately doubles and \(R^2\)
falls from 0.879 to 0.458. These protocol effects are larger than the
differences among the leading tree recipes.

The database plots explain part of this behavior. Sources are highly unequal,
conditions occupy narrow bands in oxygen--pressure space, partial-gravity
coverage is sparse, and missingness is structured by campaign. Explainability
plots identify physically plausible predictive reliance, but they cannot
separate causal mechanisms from these correlations. The defensible result is a
screening tool with explicit domain and uncertainty limits, not a universal
flammability law.
